\section{Finite results for the literature-derived formulations}
\label{sec:finite}

This section reports what is known about the question when the observation model
is a finite sample and the objective is one of the Medvedev--Brudno objects of
Section~\ref{sec:model}. The results are of two kinds. First a \emph{positive}
rigidity theorem for same-length oriented Section~6.2 candidates under
Shomorony's $\Is$. Then a sequence of \emph{negative} witnesses, each refuting a
named, precisely scoped statement. We begin with the logical device that makes
the negatives robust.

\subsection{Negative results transfer; positive results do not}

\begin{lemma}[Negative transfer]
\label{lem:negative-transfer}
Let $\Lik(\cdot;x)$ be an objective defined on a candidate class $\cands$, let
$\cands'\subseteq\cands$, and let $S,D\in\cands'$. If
$\Lik(D;x) > \Lik(S;x)$ then $S$ is not a maximizer of $\Lik(\cdot;x)$ over
$\cands$.
\end{lemma}

\begin{proof}
$D\in\cands$ is a candidate with strictly larger value than $S$, so the maximum
over $\cands$ exceeds $\Lik(S;x)$.
\end{proof}

The converse fails: a positive result over a restricted class does not imply the
same over a superclass, because the superclass may contain a better candidate.
This asymmetry is why a same-length \emph{negative} witness that is
candidate-set closed settles the unrestricted question negatively, while a
same-length \emph{positive} theorem is a restricted result unless the restriction
is itself source-justified.

\subsection{A positive same-length oriented rigidity theorem}

We work under the strict oriented single-strand convention: read types are the
oriented length-$L$ windows of $S$ with no reverse-complement collapse, and a
same-length Section~6.2 spelled candidate is a circular word $D$ of length
$G=|S|$ with $\supp\spec_L(D)=\supp\spec_L(S)$.

\begin{definition}[Window-support graph]
\label{def:graph}
Given a circular word $S$ and read length $L$, the graph $X_S$ has one node for
each distinct length-$(L-1)$ window of $S$, and one directed edge for each
distinct length-$L$ window $w\in\supp\spec_L(S)$, running from
$\operatorname{prefix}_{L-1}(w)$ to $\operatorname{suffix}_{L-1}(w)$ and carrying
multiplicity $A(w)=\spec_L(S)(w)$. A \emph{positive circulation of total $G$} is
a map $f:\supp\spec_L(S)\to\ints_{>0}$ balanced at every node with
$\sum_w f(w)=G$.
\end{definition}

The truth's spectrum $A$ is one such circulation, and the graph is strongly
connected because its edges are the edges of the single closed walk
$t\mapsto W_t(S)$.

\begin{lemma}[Same-length spelled candidates are positive circulations]
\label{lem:reduction}
Same-length Section~6.2 spelled candidates with observed support $V$ are exactly
the positive circulations of total $G$ on $X_S$.
\end{lemma}

\begin{proof}[Sketch]
A circular word $D$ of length $G$ has a balanced positive spectrum of total $G$;
if $\supp\spec_L(D)=V$ it is such a circulation. Conversely a positive
circulation on a strongly connected balanced multigraph makes it Eulerian, so an
Eulerian circuit spells a circular word $D$ of length $G$ with
$\spec_L(D)=f$. The rigidity theorem below proves uniqueness on the larger
circulation set, so it covers the spelled-circuit subclass a fortiori.
\end{proof}

\begin{theorem}[Short-repeat rigidity]
\label{thm:short-repeat}
If every length-$(L-1)$ window of $S$ occurs at most twice, then
$A=\spec_L(S)$ is the unique positive circulation of total $G$ on $X_S$.
\end{theorem}

\begin{proof}
Let $B$ be a positive circulation of total $G$ and $\delta=B-A$, a circulation
with $\sum_w\delta(w)=0$. Since $A(w)\le$ (occurrences of
$\operatorname{prefix}_{L-1}(w))\le 2$ and $B(w)\ge 1$, we have
$\delta(w)\ge-1$ on every edge.

Suppose $\delta\ne 0$. Then $N=\{w:\delta(w)<0\}$ is nonempty, and on $N$ we
have $\delta(w)=-1$, $A(w)=2$, $B(w)=1$. Take $w\in N$ with tail $v$. As
$A(w)=2$ and $\mathrm{out}_A(v)\le 2$, $w$ is the only outgoing edge at $v$, so
$\mathrm{out}_\delta(v)=-1$ and balance gives $\mathrm{in}_\delta(v)=-1$. Each
incoming edge has $\delta\ge-1$; if $v$ had two incoming edges they would each
have $A=1$ and hence $\delta\ge0$, contradicting $\mathrm{in}_\delta(v)=-1$. So
$v$ has a single incoming edge $e$ with $A(e)=2$, $\delta(e)=-1$, i.e.\ $e\in N$.
The symmetric argument at the head of $w$ shows every edge incident to a node
touched by $N$ lies in $N$; thus $N$ is a union of connected components of $X_S$
with no edge to $V\setminus N$. Strong connectivity forces $N=\emptyset$ or
$N=V$; the latter gives $\sum_w\delta(w)=-|V|<0$, contradicting
$\sum_w\delta(w)=0$. Hence $\delta\ge0$, and a nonnegative circulation with
$\sum_w\delta(w)=0$ is zero.
\end{proof}

The hypothesis of Theorem~\ref{thm:short-repeat} is exactly what the bridging
condition supplies. We record the two extension lemmas that transfer bridging
into short-repeat structure.

\begin{lemma}[Primitive extension]
\label{lem:primitive-extension}
Let $S$ be primitive and suppose some length-$(L-1)$ window occurs at least three
times. Then $S$ has a triple repeat of length at least $L-1$.
\end{lemma}

\begin{proof}[Sketch]
Choose three starts carrying the same length-$(L-1)$ window and extend the three
copies in both directions as long as all three remain equal. The process cannot
reach length $G$, since three equal length-$G$ windows would make $S$ invariant
under a nonzero shift, contradicting primitivity. At termination the preceding
symbols are not all equal and the following symbols are not all equal, which is
exactly the three-copy maximality condition; the length is at least $L-1$.
\end{proof}

\begin{lemma}[Periodic extension]
\label{lem:periodic-extension}
Let $S=P^k$ with minimal period $p<G$ and suppose some factor of $P$ of length at
least $L-1$ occurs at least twice in the circular word $P$. Then $S$ has a triple
repeat of length at least $L-1$.
\end{lemma}

\begin{proof}[Sketch]
A repeated factor of $P$ of length at least $L-1$ occurs at least four times in
$S$. Extending its two occurrence families maximally in both directions is
uniform by $p$-periodicity and cannot reach length $G$, because that would make
$r_2-r_1$ a period of $S$ and hence a smaller period of $P$. At termination the
three-copy maximality condition holds for any three occurrences meeting both
families.
\end{proof}

A length-$L$ read bridges a copy of a repeat of length $\ell$ only if
$\ell\le L-2$. Hence $R\in\Is$ forbids every triple repeat of length
$\ge L-1$. Combining this with the extension lemmas yields the main theorem.

\begin{theorem}[Rigidity under $\Is$]
\label{thm:rigidity}
If $S$ admits a read realization $R\in\Is$, then $A=\spec_L(S)$ is the unique
positive circulation of total $G$ on $X_S$.
\end{theorem}

\begin{proof}
If $S$ is primitive: $R\in\Is$ gives no triple repeat of length $\ge L-1$; by
Lemma~\ref{lem:primitive-extension} no length-$(L-1)$ window occurs three times;
Theorem~\ref{thm:short-repeat} applies. If $S=P^k$ with minimal period
$p<G$: by Lemma~\ref{lem:periodic-extension} every factor of $P$ of length
$\ge L-1$ occurs at most once, so the $p$ length-$L$ and length-$(L-1)$ windows
of $P$ are pairwise distinct and $X_S$ is a single directed $p$-cycle with
$A(w)=k$ on every edge. A positive circulation on a directed cycle is constant
by balance, and total $G=kp$ forces the value $k$, so $B=A$.
\end{proof}

\begin{corollary}[No strict same-length Section~6.2 counterexample]
\label{cor:same-length-zero}
Suppose $R\in\Is$ and the truth is Section~6.2-admissible, so
$\supp x=\supp\spec_L(S)$. Then every same-length Section~6.2 spelled candidate
$D$ satisfies $\spec_L(D)=\spec_L(S)$, and for every objective that is a function
of $(\spec_L(D),x)$,
\[
  \frac{\Lik(D;x)}{\Lik(S;x)}=1 .
\]
There is no strict same-length counterexample for any $G$, $L$, or alphabet.
\end{corollary}

\begin{proof}
By Theorem~\ref{thm:rigidity} the only positive circulation of total $G$ on
$X_S$ is $A$, so every same-length spelled candidate has spectrum $A$. Both
objectives depend on the candidate only through its spectrum, and the
observation is common, so the scores coincide.
\end{proof}

Only the triple-repeat clause of $\Is$ is used. The interleaved-repeat clause and
the coverage clause are not needed for this same-length conclusion. The equal
length assumption, by contrast, is essential and sharp; we return to it in
Section~\ref{sec:finite:length}.

\begin{proposition}[Beatability criterion]
\label{prop:beatability}
Fix $S,L$ and an observation $x$ with $\supp x=\supp\spec_L(S)$, and write
$A=\spec_L(S)$. If $A$ is the unique positive circulation of total $G$ on
$X_S$, the truth ties every same-length spelled candidate. If $B\ne A$ is
another positive circulation of total $G$, choose $w_0$ with $B(w_0)>A(w_0)$ and
put $x_M=x+M e_{w_0}$. Then the likelihood ratio equals
$\mathrm{ratio}_0\cdot(B(w_0)/A(w_0))^M$, which exceeds $1$ for all sufficiently
large $M$; the fixed-length binomial's domain is respected because $B(w_0)<G=N$
whenever $B$ is positive and non-constant.
\end{proposition}

Thus non-rigidity is exactly the same-length obstruction, and by
Theorem~\ref{thm:rigidity} the hypothesis $\Is$ removes it. The minimum
same-length non-rigid pair, found by exact search, is $S=AAAAB$,
$D=AABAB$ ($G=5$, $L=2$); its truth carries a long triple repeat and is
therefore not $\Is$-admissible. Skirting the triple-repeat clause is precisely
what permits the multiplicity reallocation.
\statusline{Theorem~\ref{thm:rigidity} and
Corollary~\ref{cor:same-length-zero} are \mathfact; the minimum non-rigid pair
and the bounded verifications are \computed, exhaustive in the cited scope.}

\subsection{Strict finite counterexamples}

The rigidity theorem is a positive same-length statement. Everything else in
this subsection is negative, and each witness is strict, hence robust to the
unresolved equivalence and tie conventions.

\begin{theorem}[Exact-multinomial counterexample over all circular candidates]
\label{thm:cex-exact}
Let $\Alphabet=\{A,B\}$, $L=3$, and let the truth be
\[
  S=AAABB \quad (G=5),
\]
with realized starts $0,1,4$, so the observed reads are $\{AAA,AAB,BAA\}$, each
once. Let
\[
  D=AAAAB \quad (|D|=5).
\]
Then $R\in\Is$ holds non-vacuously, and
\[
  \frac{\Lik_{\mathrm{ex}}(D;x)}{\Lik_{\mathrm{ex}}(S;x)} = 2 > 1 .
\]
Consequently the truth is not a maximizer for the exact multinomial over all
circular candidates, so both the maximizer and uniqueness schemas are false for
this objective.
\end{theorem}

\begin{proof}
The spectra on the observed types are $\spec_3(S)=(AAA:1,AAB:1,BAA:1)$ and
$\spec_3(D)=(AAA:2,AAB:1,BAA:1)$; $\spec_3$ has total $5$ for both. The
multinomial coefficient $\binom{3}{1,1,1}=6$ is common, and all candidates have
length $5$, so
\[
  \Lik_{\mathrm{ex}}(S;x)=6\cdot(1/5)^3=\frac{6}{125},\qquad
  \Lik_{\mathrm{ex}}(D;x)=6\cdot(2/5)(1/5)(1/5)=\frac{12}{125}.
\]
For $\Is$: the starts cover all five positions; the three $A$'s at positions
$0,1,2$ form a maximal length-$1$ triple repeat, bridged respectively by the
reads at starts $4,0,1$; and no two maximal repeat pairs have four alternating
starts, so the interleaved clause is vacuous. The all-bridged triple-repeat
clause is therefore genuinely exercised. The transfers to all circular
candidates by Lemma~\ref{lem:negative-transfer} because the witness is
same-length.
\statusline{\kernelcheck{} in the repository's
\texttt{FixedLengthExactCounterexample} module.}
\end{proof}

Repeating the start-$0$ read $k$ times gives the observed multiset
$\{AAA^k,AAB,BAA\}$ and ratio $2^k$, so the failure is a one-parameter family,
not a coincidental tie.

\begin{theorem}[Fixed-length binomial counterexample]
\label{thm:cex-binomial}
Let $\Alphabet=\{A,C\}$ inside the four-letter DNA alphabet, $L=3$, truth
$S=AAACC$ with realized starts $0,1,4$ (observed reads $AAA,AAC,CAA$), and
competitor $D=AAAAC$, both of length $5$. Under the literal Section~6.1 product
of binomial marginals with $N=5$, the zero-count factors retained,
\[
  \frac{\Lik_{\mathrm{bin}}(D;x)}{\Lik_{\mathrm{bin}}(S;x)}
  = \frac{1125}{512} > 1 .
\]
\end{theorem}

\begin{proof}[Sketch]
The instance is Theorem~\ref{thm:cex-exact} with $B$ relabeled to $C$, so the
same $\Is$ certificate transfers. The spectra are $S=(AAA:1,AAC:1,CAA:1)$ and
$D=(AAA:2,AAC:1,CAA:1)$ on the observed types. Carrying the factors for the
zero-multiplicity types changes the ratio from $2$ to $1125/512$:
\[
  \Lik_{\mathrm{bin}}(S;x)=\frac{452984832}{30517578125},\qquad
  \Lik_{\mathrm{bin}}(D;x)=\frac{7962624}{244140625}.
\]
\statusline{\kernelcheck{} in \texttt{FixedLengthBinomialCounterexample}. The
difference from the exact ratio shows that dropping or retaining the zero-count
factors changes the objective. By Lemma~\ref{lem:negative-transfer} the witness
refutes the maximizer statement over every circular candidate on which the
binomial objective is defined, since both $S$ and $D$ lie in the domain
($\spec_3(D)(AAA)=2<5$).}
\end{proof}

We now turn to the reverse-complement molecular reading of the Section~6.2 flow.
Here read types are molecule classes, the truth and competitor must both be
admissible bidirected circuits, and the objective is the literal Section~6.1
product with external $N$. The mechanism is different from the sequence-level
witnesses: reverse complementarity \emph{creates} multiplicity that the
observation does not force, and support equality then permits reallocation.

\begin{theorem}[Section~6.2 variable-length molecule counterexample]
\label{thm:cex-molecule}
Let $\Alphabet=\{A,T\}$ with involution $A\leftrightarrow T$, $L=3$, truth
$S=AAATT$ ($G=5$) with realized starts $0,1,4$, and competitor
$D=AAAATT$ ($|D|=6$). Then $R\in\Is$ holds; both $S$ and $D$ induce admissible
Section~6.2 bidirected circuits; and with external $N=5$,
\[
  \frac{\Lik_{6.1}(D)}{\Lik_{6.1}(S)} = \frac{9}{8} > 1 .
\]
\end{theorem}

\begin{proof}[Sketch]
Under $A\leftrightarrow T$ the molecule classes of $S$ are
$d_S=(AAA:1,AAT:2,TAA:2)$ and the observed classes are
$x=(AAA:1,AAT:1,TAA:1)$; for $D$ they are $d_D=(AAA:2,AAT:2,TAA:2)$. Both
spectra have the same support as $x$. The ratio of the $AAA$ factor is
$[(2/5)(3/5)^2]/[(1/5)(4/5)^2]=9/8$, and the other two classes have equal
multiplicities, so they contribute factor $1$. Bidirected-circuit feasibility and
the $\Is$ certificate are checked by exact computation and by an explicit
bidirected graph.
\statusline{\kernelcheck{} for the finite certificate and the strict inequality;
\computed{} for the bidirected graph certificate, which is checked by script and
not reproduced in Lean.}
\end{proof}

Allowing the lengths to be equal does not save the Section~6.2 reading. The
smallest same-length witness is $S=AAATAT$, $D=AAAAAT$ ($G=6$), for which the
exact ratio is $3$ and the binomial ratio is $5$, with both genomes
Section~6.2-feasible under the source's per-vertex lower bound.

\begin{remark}[A strengthening is open]
The witnesses above use the source's per-vertex lower bound, i.e.\ support
equality. Under the strictly stronger per-occurrence requirement
$d_D(w)\ge x(w)$ for every observed type, the same-length Section~6.2 statement
is \openstatus{} in the searched scope; the witnesses' truths are not candidates
under that strengthening. We flag this because it is the kind of
assumption-surface change that separates a refuted statement from an open one.
\end{remark}

\subsection{The candidate-length boundary}
\label{sec:finite:length}

The same-length restriction in the positive theorem is not cosmetic. On the next
slice, where $|D|>G$ is allowed, the truth need not be optimal even under $\Is$.

\begin{example}[Variable-length oriented boundary]
\label{ex:variable-length}
Take truth $S=AAATT$ ($G=5$), $L=3$, and competitor $D=AAAATT$ ($|D|=6$). The
supports agree and $\Is$ holds. With the observed counts
$x=\spec_3(S)+M e_{AAA}$,
\begin{center}
\begin{tabular}{@{}lll@{}}
\toprule
$M$ & exact multinomial & fixed-$N$ binomial ($N=5$) \\
\midrule
$0$ & $3125/3888<1$ & $81/128<1$ \\
$1$ & $15625/11664>1$ & $81/64>1$ \\
$2$ & $78125/34992>1$ & $81/32>1$ \\
\bottomrule
\end{tabular}
\end{center}
The boundary is exactly $n>G$: at $n=G$ the truth wins, and a single extra
observation of an over-represented type overturns it.
\statusline{\computed{}. The same word pair also appears as a
reverse-complement variable-length witness; here the orientation is fixed by
observing every truth type plus an extra $AAA$, so the mechanism is
observed-multiplicity skew rather than strand collapse.}
\end{example}

\subsection{Summary of the finite landscape}

\begin{table}[t]
\centering
\small
\begin{tabular}{@{}p{0.30\textwidth}p{0.24\textwidth}p{0.16\textwidth}p{0.20\textwidth}@{}}
\toprule
Statement & objective & candidate class & result \\
\midrule
same-length oriented Section~6.2, $R\in\Is$
  & exact / binomial
  & same-length spelled
  & \textbf{true} (Thm.~\ref{thm:rigidity}) \\
exact ML maximizer
  & exact multinomial
  & all circular
  & \textbf{false} (Thm.~\ref{thm:cex-exact}) \\
binomial ML maximizer
  & fixed-$N$ binomial
  & all circular (on its domain)
  & \textbf{false} (Thm.~\ref{thm:cex-binomial}) \\
Section~6.2 flow, molecule reads
  & Section~6.1 product
  & §6.2 circuits
  & \textbf{false} (Thm.~\ref{thm:cex-molecule}) \\
Section~6.2 flow, molecule, same length, per-vertex
  & exact / binomial
  & §6.2 spelled
  & \textbf{false} ($AAATAT\to AAAAAT$) \\
Section~6.2, per-occurrence strengthening
  & exact / binomial
  & §6.2 spelled
  & \openstatus \\
\bottomrule
\end{tabular}
\caption{The finite landscape. ``False'' means a strict counterexample over the
stated class; the evidence level (kernel-checked, exact computation, or script)
is recorded in the corresponding theorem or status line. By
Lemma~\ref{lem:negative-transfer} each same-length falsity transfers to all
circular candidates for that objective.}
\label{tab:finite}
\end{table}

Table~\ref{tab:finite} makes the shape of the answer visible. The 2016 question
is not one theorem awaiting one proof: it is a family of statements indexed by
objective, strand convention, candidate class, and length convention, some true,
some false, and at least one open. The negative results also expose the
\emph{mechanism}, which the next section isolates: the exact objective scores
empirical read-type frequencies, and a competitor can improve its score by
reallocating multiplicity toward over-represented observed types while
respecting the bridging constraint on the realized reads.
